% =========================================================================
\section{Ecuación de Laplace, cavidades y capacidad}\label{sec:cavidades}
% =========================================================================
Las propiedades anteriores son locales. Para tratar conductores con cavidades, o sistemas de varios conductores, necesitamos una herramienta
global: la formulación del problema electrostático como un problema de contorno para el potencial, y los teoremas que garantizan que su
solución es única.

% -------------------------------------------------------------------------
\subsection{Ecuaciones de Poisson y Laplace}\label{sec:poisson}
% -------------------------------------------------------------------------
\begin{proposicion}[Ecuaciones de Poisson y Laplace]\label{prop:poisson}
En electrostática, el potencial satisface la \emph{ecuación de Poisson}
\begin{equation}\label{eq:poisson}
  \resalta{\Lap\phi=-\frac{\rho}{\eo},}
\end{equation}
y, en las regiones sin carga ($\rho=0$), la \emph{ecuación de Laplace}, $\Lap\phi=0$. Las funciones de clase $C^{2}$ que satisfacen la ecuación de
Laplace en un abierto se llaman \emph{armónicas}.
\end{proposicion}
\begin{proof}
Basta sustituir $\vE=-\nab\phi$ en la ley de Gauss diferencial, $\nab\cdot\vE=\rho/\eo$:
$\nab\cdot(-\nab\phi)=-\Lap\phi=\rho/\eo$~\cite{TEMA0,TEMA1}.
\end{proof}

En presencia de conductores, el espacio vacío $\Omega=\R^{3}\setminus\bigcup_i\Vc_i$ es una región en la que (si no hay cargas libres en el vacío)
$\phi$ es armónica, y sobre cuya frontera, las superficies $\Sc_i$ de los conductores, el potencial es constante. El problema electrostático se
convierte en un \emph{problema de contorno}. Dos propiedades de las funciones armónicas son las que hacen tratable este problema.

\begin{lema}[Principio del máximo]\label{lem:maximo}
Sea $\phi$ armónica en una región acotada $\Omega$ y continua en su cierre. Entonces $\phi$ alcanza su máximo y su mínimo sobre la frontera
$\partial\Omega$, y no tiene extremos locales en el interior salvo que sea constante. Si $\Omega$ es exterior a unos conductores, el resultado se
mantiene contando el infinito como parte de la frontera.
\end{lema}
\begin{proof}
Las funciones armónicas cumplen la \emph{propiedad del valor medio}: $\phi(\vc{r}_0)$ es igual al promedio de $\phi$ sobre cualquier esfera centrada
en $\vc{r}_0$ y contenida en $\Omega$~\cite{GRIFFITHS,KELLOGG}. Si $\phi$ tuviese un máximo local estricto en $\vc{r}_0$, en una esfera pequeña alrededor de
$\vc{r}_0$ todos los valores serían menores que $\phi(\vc{r}_0)$ y su promedio también, lo que es absurdo. El caso de un máximo no estricto se reduce
al anterior por un argumento de conexión~\cite{KELLOGG}.
\end{proof}

\begin{teorema}[Teoremas de unicidad]\label{teo:unicidad}
Sea $\Omega$ la región exterior a un conjunto de conductores $\Vc_1,\dots,\Vc_N$, con una densidad de carga $\rho$ dada en $\Omega$ y la condición
$\phi\to0$ en el infinito. La solución de la ecuación de Poisson en $\Omega$ es única si se conoce, para cada conductor:
\begin{enumerate}[label=\textup{(\alph*)}]
  \item su potencial $\phi_i$ (\emph{primer teorema de unicidad}, problema de Dirichlet), o bien
  \item su carga total $Q_i$ (\emph{segundo teorema de unicidad}),
\end{enumerate}
o una combinación de ambos datos (el potencial para unos conductores y la carga para los demás).
\end{teorema}
\begin{proof}
Sean $\phi_a$ y $\phi_b$ dos soluciones y $\psi=\phi_a-\phi_b$. Entonces $\Lap\psi=0$ en $\Omega$, $\psi\to0$ en el infinito y $\psi$ es constante, de
valor $\psi_i$, sobre cada $\Sc_i$. Aplicamos el teorema de la divergencia al campo $\psi\nab\psi$, cuya divergencia es
$\nab\cdot(\psi\nab\psi)=\norm{\nab\psi}^{2}+\psi\Lap\psi=\norm{\nab\psi}^{2}$, en la región comprendida entre los conductores y una esfera de radio
$R\to\infty$ (en la que el integrando decae como $R^{-3}$ y su contribución se anula):
\[
  \iiint_{\Omega}\norm{\nab\psi}^{2}\dV=-\sum_{i=1}^{N}\oiint_{\Sc_i}\psi\,\nab\psi\cdot\un\,\dsc
  =-\sum_{i=1}^{N}\psi_i\oiint_{\Sc_i}\nab\psi\cdot\un\,\dsc ,
\]
donde $\un$ es la normal exterior a cada conductor; el signo menos aparece porque la normal exterior a la región $\Omega$ sobre $\Sc_i$ es
$-\un$. Ahora bien: (a) si los potenciales coinciden, $\psi_i=0$; (b) si las cargas coinciden, como $Q_i=\eo\oiint_{\Sc_i}\vE\cdot\un\,\dsc
=-\eo\oiint_{\Sc_i}\nab\phi\cdot\un\,\dsc$ por \cref{eq:campo-superficie}, se tiene
$\oiint_{\Sc_i}\nab\psi\cdot\un\,\dsc=-\frac{1}{\eo}\bigl(Q_i^{(a)}-Q_i^{(b)}\bigr)=0$. En cualquiera de los dos casos, y en las combinaciones, cada sumando se
anula, de modo que $\int_{\Omega}\norm{\nab\psi}^{2}=0$, $\nab\psi=\vc{0}$ y $\psi$ es constante; como $\psi\to0$ en el infinito, $\psi\equiv0$.
\end{proof}

\begin{observacion}[El significado práctico de la unicidad]\label{obs:unicidad}
La unicidad legitima el método de «adivinar y comprobar»: cualquier función que cumpla la ecuación y las condiciones de contorno es \emph{la}
solución. Es la base del método de las imágenes, que usaremos en la \cref{fig:corteza}, y del principio de superposición para conductores,
que da lugar a los coeficientes de potencial (\cref{sec:sistemas}).
\end{observacion}

% -------------------------------------------------------------------------
\subsection{Conductores con cavidades}\label{sec:cavidades-casos}
% -------------------------------------------------------------------------
Consideramos un conductor con una cavidad interior (una región vacía completamente rodeada de metal). Su superficie tiene dos partes: la
\emph{cara exterior} y la \emph{cara interior}, que limita la cavidad.

\begin{teorema}[Cavidad vacía: jaula de Faraday]\label{teo:faraday}
Si una cavidad de un conductor en equilibrio no contiene carga, el campo en su interior es nulo y la cara interior no tiene carga en ningún
punto, cualesquiera que sean las cargas exteriores y la carga del conductor:
\begin{equation}\label{eq:faraday}
  \resalta{\Ecav=\vc{0},\qquad\sigma\big|_{\text{cara interior}}=0 .}
\end{equation}
\end{teorema}
\begin{proof}
En la cavidad $\phi$ es armónica, y en su frontera (la cara interior) vale la constante $\phi_{\Vc}$, el potencial del conductor. La función
constante $\phi\equiv\phi_{\Vc}$ cumple la ecuación y la condición de contorno; por el primer teorema de unicidad (o, directamente, por el principio
del máximo: el máximo y el mínimo en la cavidad se alcanzan en la frontera, donde valen lo mismo) es la única solución. Por tanto
$\Ecav=-\nab\phi=\vc{0}$, y por \cref{eq:campo-superficie} $\sigma=\eo\,\vE\cdot\un=0$ en la cara interior.

\emph{Argumento alternativo con líneas de campo.} Si hubiese campo en la cavidad, sus líneas tendrían que nacer y morir en la cara interior
(no hay cargas dentro de la cavidad), es decir, saldrían del conductor y volverían a él, lo que prohíbe el \cref{cor:lineas}.
\end{proof}

\begin{proposicion}[Regla de las caras]\label{prop:caras}
Sea un conductor de carga total $Q$ con una cavidad que contiene una carga total $q$. Entonces, en el equilibrio,
\begin{equation}\label{eq:caras}
  \resalta{\Qint=-q,\qquad\Qext=Q+q .}
\end{equation}
\end{proposicion}
\begin{proof}
Tomamos una superficie gaussiana $S_G$ contenida por completo en el metal y que rodee la cavidad (curvas a trazos de la \cref{fig:cavidades}).
Sobre ella $\vE=\vc{0}$, luego $\Qenc=\eo\oiint_{S_G}\vE\cdot\dS=0$. La carga encerrada es la de la cavidad más la de la cara interior, $q+\Qint=0$, de donde
$\Qint=-q$. Como el metal no tiene carga en volumen (\cref{teo:propiedades}), la carga del conductor se reparte entre sus dos caras,
$Q=\Qint+\Qext$, y $\Qext=Q-\Qint=Q+q$.
\end{proof}

% =========================================================================
%  Figura: los cuatro casos de conductor hueco (como en las figuras del tema).
%  Contornos en polares (centro común):
%    cara exterior  r(t) = 1.50 + 0.16 cos 3t + 0.10 sen 2t
%    gaussiana      r(t) = 1.06 + 0.08 cos 3t + 0.05 sen 2t
%    cara interior  r(t) = 0.62 + 0.08 cos(2t + 40)
% =========================================================================
\newcommand{\cavExt}{plot[domain=0:360, samples=90, smooth cycle, variable=\t]
  ({(1.50+0.16*cos(3*\t)+0.10*sin(2*\t))*cos(\t)},{(1.50+0.16*cos(3*\t)+0.10*sin(2*\t))*sin(\t)})}
\newcommand{\cavInt}{plot[domain=0:360, samples=70, smooth cycle, variable=\t]
  ({(0.62+0.08*cos(2*\t+40))*cos(\t)},{(0.62+0.08*cos(2*\t+40))*sin(\t)})}
\newcommand{\cavGauss}{plot[domain=0:360, samples=80, smooth cycle, variable=\t]
  ({(1.06+0.08*cos(3*\t)+0.05*sin(2*\t))*cos(\t)},{(1.06+0.08*cos(3*\t)+0.05*sin(2*\t))*sin(\t)})}
% signos sobre la cara exterior (#1 = 1) o interior (#1 = 2), en el lado del metal
\newcommand{\signosExt}{\foreach \t in {0,30,...,330}{
  \node[signopos, font=\tiny\bfseries] at ({(1.36+0.16*cos(3*\t)+0.10*sin(2*\t))*cos(\t)},{(1.36+0.16*cos(3*\t)+0.10*sin(2*\t))*sin(\t)}) {$+$};}}
\newcommand{\signosInt}{\foreach \t in {0,36,...,324}{
  \node[signoneg, font=\tiny\bfseries] at ({(0.76+0.08*cos(2*\t+40))*cos(\t)},{(0.76+0.08*cos(2*\t+40))*sin(\t)}) {$-$};}}
\newcommand{\panelCavidad}[4]{% #1 título, #2 signos ext (0/1), #3 carga en cavidad (0/1), #4 rótulos
  \fill[metal!55!white, even odd rule] \cavExt \cavInt;
  \fill[pattern=north east lines, pattern color=azulcielo!70, even odd rule] \cavExt \cavInt;
  \draw[azulprof, line width=0.6pt] \cavExt;
  \draw[azulprof, line width=0.6pt] \cavInt;
  \draw[gauss] \cavGauss;
  \ifnum#2=1 \signosExt\fi
  \ifnum#3=1 \signosInt \node[qpos, minimum size=8pt] at (0.05,-0.02) {$+$};
             \node[font=\scriptsize, text=cargapos] at (0.05,0.3) {$q$};\fi
  \node[font=\scriptsize, text=azulprof, fill=white, inner sep=1pt, rounded corners=1pt] at (0.15,-1.25) {$\vE=\vc{0}$};
  \node[anchor=west, align=left, font=\footnotesize] at (1.95,0) {#4};
  \node[anchor=south, font=\small] at (1.15,1.85) {#1};}

\begin{figure}[t]
\centering
\begin{tikzpicture}[x=0.95cm, y=0.95cm]
  \begin{scope}
    \panelCavidad{(a) Neutro, cavidad vacía}{0}{0}
      {$\Qint=0$\\ $\Qext=0$\\[2pt] $\Ecav=\vc{0}$\\ $\Eext=\vc{0}$}
  \end{scope}
  \begin{scope}[xshift=8.2cm]
    \panelCavidad{(b) Carga $Q$, cavidad vacía}{1}{0}
      {$\Qint=0$\\ $\Qext=Q$\\[2pt] $\Ecav=\vc{0}$\\ $\Eext\neq\vc{0}$}
  \end{scope}
  \begin{scope}[yshift=-4.4cm]
    \panelCavidad{(c) Neutro, carga $q$ en la cavidad}{1}{1}
      {$\Qint=-q$\\ $\Qext=+q$\\[2pt] $\Ecav\neq\vc{0}$\\ $\Eext\neq\vc{0}$}
  \end{scope}
  \begin{scope}[xshift=8.2cm, yshift=-4.4cm]
    \panelCavidad{(d) Carga $Q$, carga $q$ en la cavidad}{1}{1}
      {$\Qint=-q$\\ $\Qext=Q+q$\\[2pt] $\Ecav\neq\vc{0}$\\ $\Eext\neq\vc{0}$}
  \end{scope}
  % leyenda común
  \draw[gauss] (2.6,-7.05) -- (3.4,-7.05);
  \node[anchor=west, font=\footnotesize] at (3.5,-7.05) {superficie de Gauss (cerrada) dentro del metal};
\end{tikzpicture}
\caption{Conductores huecos en equilibrio. La superficie de Gauss trazada dentro del metal, donde $\vE=\vc{0}$, encierra carga nula: la cara
interior lleva siempre la carga opuesta a la de la cavidad, y la exterior el resto de la carga del conductor. Una cavidad vacía no tiene campo.}
\label{fig:cavidades}
\end{figure}


La \cref{tab:cavidades} resume los cuatro casos posibles, que la \cref{fig:cavidades} representa.

\begin{table}[ht]
\centering
\caption{Carga de cada cara y campo en la cavidad y en el exterior de un conductor hueco aislado, sin más cargas que las indicadas.}
\label{tab:cavidades}
\begin{tabular}{@{}lcccc@{}}
  \toprule
  Situación & Cara interior & Cara exterior & $\Ecav$ & $\Eext$\\
  \midrule
  Neutro, cavidad vacía                 & $0$  & $0$    & $\vc{0}$        & $\vc{0}$\\
  Carga $Q$, cavidad vacía              & $0$  & $Q$    & $\vc{0}$        & $\neq\vc{0}$\\
  Neutro, carga $q$ en la cavidad       & $-q$ & $+q$   & $\neq\vc{0}$    & $\neq\vc{0}$\\
  Carga $Q$, carga $q$ en la cavidad    & $-q$ & $Q+q$  & $\neq\vc{0}$    & $\neq\vc{0}$\\
  \bottomrule
\end{tabular}
\end{table}

\begin{teorema}[Apantallamiento en los dos sentidos]\label{teo:apantallamiento}
Sea un conductor hueco cerrado, con cargas en su cavidad y en el exterior.
\begin{enumerate}[label=\textup{(\alph*)}]
  \item El campo en la cavidad depende solo de las cargas de la cavidad y de la forma de la cara interior: no se ve afectado por las cargas
        exteriores (salvo por una constante aditiva en el potencial).
  \item El campo exterior depende solo de las cargas exteriores, de la forma de la cara exterior y de la carga total $\Qext=Q+q$ de esta cara: no
        depende de cómo estén colocadas las cargas dentro de la cavidad.
\end{enumerate}
\end{teorema}
\begin{proof}
(a) En la cavidad, $\phi-\phi_{\Vc}$ cumple la ecuación de Poisson con las cargas de la cavidad y se anula sobre la cara interior: por el primer
teorema de unicidad está determinada sin referencia alguna al exterior. (b) En el exterior, $\phi$ cumple la ecuación de Poisson con las cargas
exteriores, tiende a cero en el infinito, es constante sobre la cara exterior y la carga de esa cara es $\Qext=Q+q$ (\cref{prop:caras}): por el segundo
teorema de unicidad el problema exterior queda determinado, y en sus datos solo interviene el valor $q$, no la posición de la carga.
\end{proof}

La \cref{fig:corteza} ilustra el teorema con una corteza esférica gruesa, que puede resolverse exactamente por el método de las imágenes~\cite{JACKSON}.
Con una carga $q$ en el centro (a) todo tiene simetría esférica. Al desplazar la carga (b), las líneas de la cavidad se deforman, pero fuera de la
corteza el campo sigue siendo exactamente el de una carga $q$ en el centro: la cara exterior «no se entera». Al añadir una carga exterior (c), el
campo exterior se deforma, pero el de la cavidad es idéntico al de (b). Por último, sin carga en la cavidad (d), el interior queda completamente
apantallado: es la jaula de Faraday.

% =========================================================================
%  Figura «tipo Mathematica»: corteza esférica conductora gruesa (radios
%  R1 = 1 de la cavidad y R2 = 2 exterior), neutra, en el plano meridiano.
%  Solución exacta por imágenes:
%   - cavidad: q en s u_x  +  imagen -q R1/s en (R1^2/s) u_x;
%   - exterior: Qe en -D u_x, imagen -Qe R2/D en -(R2^2/D) u_x y la carga
%     q + Qe R2/D en el centro (la corteza aislada tiene Q_ext = q).
%  Datos: figuras/datos/corteza_{a,b,c,d}_*.
% =========================================================================
\newcommand{\panelCorteza}[2]{% #1 = a, b, c, d ; #2 = título
  \fill[azulfondo!60!white] (-5.4,-3.3) rectangle (3.4,3.3);
  \begin{scope}
    \clip (-5.4,-3.3) rectangle (3.4,3.3);
    \input{figuras/datos/corteza_#1_lc}
    \input{figuras/datos/corteza_#1_fl}
  \end{scope}
  \fill[metal!70!white, even odd rule] (0,0) circle (2) (0,0) circle (1);
  \fill[top color=azulvelo!30, bottom color=metalosc!45, opacity=0.6, even odd rule] (0,0) circle (2) (0,0) circle (1);
  \draw[azulprof, line width=0.55pt] (0,0) circle (2) (0,0) circle (1);
  \input{figuras/datos/corteza_#1_q}
  \draw[azulclaro, line width=0.5pt] (-5.4,-3.3) rectangle (3.4,3.3);
  \node[anchor=north, font=\small] at (-1.0,-3.4) {#2};}

\begin{figure}[t]
\centering
\begin{tikzpicture}[x=0.74cm, y=0.74cm]
  \tikzset{numLC/.style={azulmedio, line width=0.42pt},
           numFL/.style={-{Stealth[length=3pt, width=2.5pt]}, azulmedio, line width=0.42pt}}
  \begin{scope}
    \panelCorteza{a}{(a) Carga $q$ en el centro}
    \node[qpos, minimum size=7pt, font=\tiny\bfseries] at (0,0) {$+$};
  \end{scope}
  \begin{scope}[xshift=8.9cm]
    \panelCorteza{b}{(b) Carga $q$ desplazada}
    \node[qpos, minimum size=7pt, font=\tiny\bfseries] at (0.45,0) {$+$};
  \end{scope}
  \begin{scope}[yshift=-6.0cm]
    \panelCorteza{c}{(c) Carga $q$ desplazada y carga exterior $Q_e$}
    \node[qpos, minimum size=7pt, font=\tiny\bfseries] at (0.45,0) {$+$};
    \node[qpos, minimum size=8pt, font=\tiny\bfseries] at (-4,0) {$+$};
    \node[fondo, font=\scriptsize, text=cargapos] at (-4,0.52) {$Q_e$};
  \end{scope}
  \begin{scope}[xshift=8.9cm, yshift=-6.0cm]
    \panelCorteza{d}{(d) Cavidad vacía y carga exterior $Q_e$}
    \node[qpos, minimum size=8pt, font=\tiny\bfseries] at (-4,0) {$+$};
    \node[fondo, font=\scriptsize, text=cargapos] at (-4,0.52) {$Q_e$};
    \node[font=\scriptsize, text=azulprof, fill=white, inner sep=1pt, rounded corners=1pt] at (0,0) {$\Ecav=\vc{0}$};
  \end{scope}
\end{tikzpicture}
\caption{Corteza esférica conductora neutra (radios interior $R_1$ y exterior $R_2=2R_1$), resuelta exactamente por el método de las imágenes.
(a)~Carga $q$ en el centro: simetría esférica. (b)~Al desplazar $q$ cambia el campo de la cavidad, pero el exterior sigue siendo el de una carga
$q$ en el centro. (c)~Una carga exterior $Q_e$ deforma el campo exterior, pero no altera el de la cavidad. (d)~Sin carga en la cavidad, el interior
queda apantallado (jaula de Faraday). Se indican los signos de la carga inducida en cada cara.}
\label{fig:corteza}
\end{figure}


\begin{observacion}[Aplicaciones del apantallamiento]\label{obs:jaula}
Una caja metálica cerrada protege su interior de cualquier campo electrostático exterior: es el principio de las jaulas de Faraday, de los cables
coaxiales con malla exterior y de la protección de los ocupantes de un coche o un avión alcanzado por un rayo. En la práctica, una malla con aberturas
pequeñas frente a la escala de variación del campo funciona casi igual que una superficie continua. Recíprocamente, conectando el conductor a tierra
($\phi_{\Vc}=0$, $\Qext=0$) también se impide que las cargas de la cavidad creen campo fuera.
\end{observacion}

% -------------------------------------------------------------------------
\subsection{Capacidad de un conductor aislado}\label{sec:capacidad}
% -------------------------------------------------------------------------
Sea un conductor aislado, cargado con carga $Q$, en equilibrio. El potencial que crea en un punto exterior es, por la ley de Coulomb para
distribuciones superficiales~\cite{TEMA1},
\begin{equation}\label{eq:pot-conductor}
  \phi(\vc{r})=\frac{1}{4\pi\eo}\oiint_{\Sc}\frac{\sigma(\rp)}{\norm{\vc{r}-\rp}}\,\dSp ,
\end{equation}
y, si $\vc{r}\in\Sc$, se obtiene el potencial del propio conductor, $\phi_{\Vc}$.

\begin{teorema}[Capacidad]\label{teo:capacidad}
Para un conductor aislado, el cociente entre su carga y su potencial no depende de $Q$:
\begin{equation}\label{eq:capacidad}
  \resalta{C=\frac{Q}{\phi_{\Vc}}=\frac{Q'}{\phi'_{\Vc}}=\mathrm{cte},}
\end{equation}
y solo depende de la geometría del conductor (y del medio). Se llama \emph{capacidad} del conductor; su unidad en el SI es el
\emph{faradio}, $1\,\mathrm{F}=1\,\mathrm{C/V}$.
\end{teorema}
\begin{proof}
Sea $\sigma$ la densidad de equilibrio cuando la carga es $Q$, con potencial $\phi_{\Vc}$. Si la carga pasa a ser $Q'=\lambda Q$, la densidad
$\lambda\sigma$ tiene carga total $\lambda Q$ y, por la linealidad de \cref{eq:pot-conductor}, crea un potencial $\lambda\phi$, que es constante
($\lambda\phi_{\Vc}$) sobre el conductor y nulo en el infinito. Por el segundo teorema de unicidad es la solución de equilibrio: la geometría no cambia,
luego la forma en que se distribuye la carga tampoco, y solo se reescala. Así $\phi'_{\Vc}=\lambda\phi_{\Vc}$ y $Q'/\phi'_{\Vc}=Q/\phi_{\Vc}$.
\end{proof}

\begin{ejemplo}[Esfera conductora aislada]\label{ej:esfera}
Para una esfera de radio $R$, la carga se reparte uniformemente y el potencial es el de una carga puntual en el centro evaluado en $r=R$,
$\phi=\kc Q/R$. Así
\begin{equation}\label{eq:cap-esfera}
  C_{\text{esfera}}=4\pi\eo R .
\end{equation}
El faradio es una unidad enorme: la propia Tierra, tratada como una esfera conductora de radio $6{,}37\cdot10^{6}\,\mathrm{m}$, tiene una capacidad de solo
$C\simeq7{,}1\cdot10^{-4}\,\mathrm{F}\simeq709\,\mu\mathrm{F}$. En la práctica se usan sus submúltiplos: $1\,\mu\mathrm{F}=10^{-6}\,\mathrm{F}$,
$1\,\mathrm{nF}=10^{-9}\,\mathrm{F}$ y $1\,\mathrm{pF}=10^{-12}\,\mathrm{F}$.
\end{ejemplo}
